We provide a table listing all notations and their explanations in Table~\ref{tab:notations}. In Section~\ref{sec:appendix:kernel_comparison}, we compare~\Algname's kernel design with other open-source MoE kernel designs. In Section~\ref{sec:appendix:dH_proof}, we elaborate on \Algname's computational path for $dS$ and $dH$ that does not use $Y$ and $dY$. In Section~\ref{sec:appendix:dS_computation}, we justify \Algname's computational path for $dS$ is both activation memory and computationally-efficient. In Section~\ref{sec:appendix:top-k}, We examine \Algname's top-$K$ sorting kernel. In Table~\ref{tab:moe-scaling}, we provide a trending overview for open-source frontier MoE models. We present \Algname's expert aggregation strategy in Figure~\ref{fig:expert-agg}. Figure~\ref{fig:expert-agg-why-slow} illustrates that on Hopper GPUs, asynchronous TMA store (top) has higher memory bandwidth and can naturally overlap with TensorCore MMA. In addition, \Algname's up-projection backward is included in Algorithm~\ref{alg:up-bwd-moe}. In Section~\ref{sec:appendix:kernel_ablation}, we present ablation studies of training throughput for \Algname's MoE computation kernels to examine the impact of each design choice made for~\Algname. In Section~\ref{sec:appendix:more_experiments}, we assess the quality improvements of MoE models trained by varying expert granularity. We then focus on various ablation studies on our token rounding routing algorithm to assess the quality difference of the trained MoE models from the choice of rounding subroutine. We also study the effect of microbatch size $T$ and tile size $\tileM$ on token rounding. In Section~\ref{sec:appendix:config}, we describe the configurations for benchmarking the memory and training throughput. In Section~\ref{sec:appendix:exp-details}, we include the details of model training and the evaluation setup. 

% We provide an asymptotic scaling of wasted FLOPs by tile quantization effect in Section~\ref{sec:appendix:wasted_flops}.